\scanpage{72}

thus \unsure{an} exact sequence of $\mathbb Z$-modules
\diagram{\begin{tikzcd}[column sep=small,row sep=small]
{(131)\quad1}\arrow[r,dash]
&\mathbb Z\arrow[r]
&N^{0+}\arrow[r]
&\Gamma^{0+}\arrow[r]\arrow[d,phantom,"{\rotatebox{90}{$\simeq$}\ \mathrm{dfn}}" description]
&0\\
&n\arrow[r,mapsto]&(n\lambda_0,n\lambda_0)
&N^{0+}/\operatorname{Im}\mathbb Z&
\end{tikzcd}}
with $\Gamma^{0+}$ a free $\mathbb Z$-module of rank $1$.

We conclude from this
\[
(132)\qquad\Gamma_{\rho,\sigma}^{0+}\simeq\underline\Gamma^{0+}
\qquad(\simeq\mathbb G_a)
\]

\deleted{\unsure{Hence}} Thus, all the group schemes we have just constructed are smooth over $\operatorname{Spec}\mathbb Z$.

Can we define a more-or-less natural section of the extension
(\inserted{$N_{\rho,\sigma}$ of $\Gamma_{\rho,\sigma}$ by $L_0^{\underline M}$, hence of}
$M_{\rho,\sigma}^{\sharp}$ of
$\underbrace{\Gamma_{\rho,\sigma}}_{\inserted{\substack{\text{smooth of relative}\\\text{dimension }2}}}$
by $\underline M$?

Taking into account the loop equation which describes $N_{\rho,\sigma}$, it seems that we \unsure{find} such a section \unsure{as soon|some} as we choose a splitting of
\diagram{\begin{tikzcd}[column sep=small,row sep=large]
{} &0\arrow[r]
&\overset{\inserted{\mathrm{rg}\,2}}{N^{0+}}\arrow[r]
&M\times M\arrow[r]
&\overset{\inserted{\mathrm{rg}\,2}}{M_0}\arrow[r]
&0&\text{and of}\\
{(133)}\arrow[u,phantom,"\Bigg\{" description]
&0\arrow[r]
&\underset{\inserted{\mathrm{rg}\,1}}{\mathbb Z}\arrow[r]
&N^{0+}\arrow[r]
&\underset{\inserted{\mathrm{rg}\,1}}{\Gamma^{0+}}\arrow[r]
&0&
\end{tikzcd}}
(category of $\mathbb Z$-modules). But we must check that the \textquotedblleft set-theoretic\textquotedblright{} sections of schemes that \unsure{this} allows us to define are multiplicative\dots

\scanpage{73}
\mspage{684}

\(6^\circ)\) \(\underline Z_{\rho,\sigma}\) and
\(\underline G_{\rho,\sigma}=\mathbb O_{12}^{*}\mathbin{.}\underline{\mathcal G}^{+},\quad
\underline G_{\rho,\sigma}^{\sharp}\simeq\mathbb O_{12}^{*}\mathbin{.}\underline{\mathcal G}^{\sharp}
=\mathbb O_{12}^{*}\mathbin{.}\underline M\).

Once \(\mathbb O_{12}^{*}\) and \(M_0\subset M\) have been introduced, it
seems to me that the group \(\underline Z_{\rho,\sigma}\) (in place
of \(\underline{\widetilde Z}_{\rho,\sigma}\), in which we took
account of the \(\alpha,\beta\) and not only
of \(\xi,\eta\)) \unsure{again becomes} nice:
\[
(134)\quad
\begin{gathered}
\underline Z_{\rho,\sigma}\subset
\underline{\operatorname{Aut}}(\underline{\mathcal G}^{+})\times\mathbb O_{12}^{*}
\\[-2pt]
\text{\(\prime\prime\)}\\[-2pt]
\left\{(u,\underline\mu)\ \left|\
\begin{array}{ll}
\text{a)}&u(\rho)=\xi\rho\quad\text{(locally) }\underline{\mathcal G}_{\tau}^{\sharp}\text{-conjugate to }\rho,\text{ by an}\\
&\alpha\in\underline{\mathcal G}_{\tau}^{\sharp}\text{ such that }\unsure{\chi}(\alpha)=\varepsilon_3(\underline\mu)\\
\text{b)}&u(\sigma)=\eta\sigma\quad\text{(locally) }\underline{\mathcal G}_{\tau}^{\sharp}\unsure{+|-}\text{conjugate to }\sigma,\text{ by a}\\
&\beta\in\underline{\mathcal G}_{\tau}^{\sharp}\text{ such that }\chi(\beta)=\varepsilon_2(\underline\mu)\\
\text{c)}&\text{loop equation}\\
&\xi_0=\eta_0(\lambda_{\unsure{1}}^3)^{\nu''}\quad\text{i.e.}\quad\xi_0=\eta_0+3\nu''\lambda_1\\
&\hfill\text{in }\underline M_0
\end{array}\right.\right.
\end{gathered}
\]
where \(\nu''=\nu''(\underline\mu)\in\Gamma\mathbb G_a\)
(cf. \unsure{page\ verso} page~686), and \unsure{above\(^{t}\)})
\[
(135)\qquad
\left\{\begin{aligned}
\xi_0&=\xi-\nu_3\lambda_1\in\Gamma\underline M_0
&\nu_3&=\nu_3(\underline\mu)\\
\eta_0&=\eta\omega^{\nu_2}\in\Gamma\underline M_0
&\nu_2&=\nu_2(\underline\mu)
\end{aligned}\right.
\]
\deleted{We have that} \inserted{We shall have} \deleted{must suppose} moreover that
\[
(136)\qquad
\left\{\begin{aligned}
&\xi_0\quad\text{\unsure{locally|locally\(^{mt}\)} in the image of }(1-\rho_{\underline M})(\underline M)\\
&\eta_0\quad\text{\unsure{locally|locally\(^{mt}\)} in the image }(1-\sigma_{\underline M})(\underline M)
\end{aligned}\right.
\]
and we shall have
\[
(137)\quad
\begin{aligned}
\underline Z_{\rho,\sigma}\simeq\underline\Sigma_{\rho,\sigma}^{\mathrm{eff}}
\overset{\mathrm{dfn}}{=}
\bigl\{(\underline\mu,\xi_0,\eta_0)&\in\Gamma\mathbb O_{12}^{*}\times\underline M_0\times M_0\ \big|\\
1^\circ)\quad\xi_0&=\eta_0+\deleted{\unsure{\xi}}\nu''(3\lambda_1)\quad\text{in }\underline M_0\\
&\text{where }\nu''=\nu''(\underline\mu)\\
2^\circ)\quad&\text{we have }(136)\quad\deleted{\unsure{\bigr\}}}\ \bigr\}.
\end{aligned}
\]
Now

\((1-\rho_{\underline M}):\)
\diagram{\begin{tikzcd}[column sep=large]
\underline M\arrow[r]&\underline M_0
\end{tikzcd}}

\((1-\sigma_{\underline M}):\)
\diagram{\begin{tikzcd}[column sep=large]
\underline M\arrow[r]&\underline{\mathbb Z}=\mathbb G_a\arrow[r]&\underline M_0\\
(\unsure{b_0}\lambda_0+b_1\lambda_1)\arrow[r]&(b_1-\underset{\deleted{\unsure{4|\lambda}}}{b_0})\,\deleted{\unsure{\lambda}}\arrow[r,mapsto]&\lambda\underbrace{(\lambda_1-\lambda_0)}_{\delta}=\lambda\delta
\end{tikzcd}}

\scanpage{74}

thus (136) \deleted{\unsure{is written}} is a condition on
\(\eta_0\) alone, namely

\((136\mathrm{bis})\quad\eta_0\in\Gamma\operatorname{Im}\bigl(\)
\diagram{\begin{tikzcd}[column sep=large]
\mathbb G_a\arrow[r]&\underline M_0\\
\lambda\arrow[r,mapsto]&\lambda\delta
\end{tikzcd}}
\(\bigr)=\unsure{\Gamma}\operatorname{Ker}\bigl(\)
\diagram{\begin{tikzcd}[column sep=large]
M_0\arrow[r]&\mathbb G_a\\
c_0\lambda_0+c_1\lambda_1\arrow[r]&c_0+c_1
\end{tikzcd}}
\(\bigr)\)

\deleted{That said, we thus find \unsure{as}}
\inserted{Let us note that we have}
\[
\deleted{(138)\qquad\underline\Sigma_{\rho,\sigma}^{\mathrm{eff}}\simeq}
\]
\[
(138)\qquad
\left\{\begin{aligned}
M_0&=\mathbb Z\underbrace{\delta}_{\lambda_1-\lambda_0}
\oplus\mathbb Z\mathbin{.}\underbrace{(3\lambda_1)}_{(1-\rho)\delta}
\qquad\text{whence}\\
\underline M_0&=\deleted{\unsure{\mathbb G_a}}\ \underline{\mathbb Z\delta}
\times\underline{\mathbb Z\mathbin{.}(3\lambda_1)}\ .
\end{aligned}\right.
\]
and \deleted{consequently} we find

\((139)\)
\diagram{\begin{tikzcd}[column sep=large]
\underline\Sigma_{\rho,\sigma}^{\mathrm{eff}}\arrow[r,"\simeq"]&
\mathbb O_{12}^{*}\times\deleted{\unsure{\underline{\mathbb Z}\delta}}\ \deleted{\unsure{\mathbb G_a}}\ \underset{\unsure{\in|\simeq}\mathbb G_a}{\underline{\mathbb Z\delta}}\\
(\underline\mu,\xi_0,\unsure{\eta_0|\beta_0})\arrow[r,mapsto]&(\underline\mu,\unsure{\eta_0|\beta_0})
\end{tikzcd}}
.

For \(\unsure{\xi_0|\zeta_0}\) is recovered with \(\underline\mu,\unsure{\eta_0|\beta_0}\)
\deleted{\unsure{from}} the loop equation
\[
\deleted{\boxed{\begin{aligned}
\unsure{\alpha_0|\xi_0}&=(1-\rho_{\underline M})^{-1}
\bigl(\unsure{\eta_0|\beta_0}+\nu''\underbrace{(3\lambda_1)}_{(1-\rho_{\underline M})(\delta)}\bigr)\\
&\deleted{\unsure{=\nu''\delta+(1-\rho_{\underline M})^{-1}\eta_0}}\\
&=\nu''\delta+\unsure{\xi_0|\eta_0}
\end{aligned}}}
\]

\uline{N.B.} If on the other hand we give ourselves \(\xi_0\), we obtain from it
by the loop equation and (138) the value of \(\unsure{\eta_0|\beta_0}\),
and that of \(\nu''\) -- thus if we know
in addition \(\nu_2\) and \(\nu_3\), we deduce
the value of \(2\nu_0\) hence (over a
base where 2 is a regular section of \(\mathscr O_S\)) \(\nu_0\),
i.e. finally \(\underline\mu\). Thus over a

\scanpage{75}
\mspage{685}

base \(S\) where 2 is regular, we find

\((140)\)
\diagram{\begin{tikzcd}[column sep=large]
\underline\Sigma_{\rho,\sigma}^{\mathrm{eff}}(S)\arrow[r,hook]&
\bigl(\underline M_0\deleted{\unsure{(S)}}\times\{0,1\}\deleted{\unsure{(S)}}\times\{0,1\}\bigr)(S)\\
(\underline\mu,\xi_0,\eta_0)\arrow[r,mapsto]&(\xi_0,\varepsilon_2(\underline\mu),\varepsilon_3(\underline\mu))
\end{tikzcd}}

the image being formed of the systems
\((\xi_0,\unsure{\varepsilon_2|\xi_2},\unsure{\varepsilon_3|\xi_3})\) such that, setting
\[
\deleted{(\unsure{141|140})}\qquad
\xi_0=\unsure{\eta_0|\beta_0}\delta+\nu''(3\lambda_1)
\qquad\eta_0,\nu''\in\deleted{\unsure{\Gamma}}\mathbb G_a(S)
\]
thanks to (138), we have
\[
(141)\qquad
\left\{\begin{aligned}
&\nu''+\nu_2-\nu_3[+2\nu_2\nu_3]\quad
\begin{gathered}
\inserted{\text{divisible by 2}}\\[-2pt]
\deleted{\text{of the form }2\nu_0}\\[-2pt]
\deleted{\text{divisible by}}
\end{gathered}\\
&\underbrace{(-5)^{\nu_2}(+5)^{\nu_3}+6(\nu''+\nu_2-\nu_3+2\nu_2\nu_3)}_{1-2\nu_3-2\nu_2\nu_3+6\nu''}\quad\text{inverse}
\end{aligned}\right.
\]

In any case we see that
\(\underline\Sigma_{\rho,\sigma}^{\mathrm{eff}}\) hence also \(\underline Z_{\rho,\sigma}\) is smooth of
relative dimension 2

\((142)\)
\diagram{\begin{tikzcd}[column sep=large]
1\arrow[r,dash]&\unsure{\underline Z}_{\rho,\sigma}^{+0}\arrow[r]\arrow[d,phantom,"\rotatebox{90}{$\simeq$}"']&
\underline Z_{\rho,\sigma}\arrow[r]&\mathbb O_{12}^{\unsure{*|\circ}}\arrow[r,dash]&1\\
&\mathbb G_a\simeq\underline{\mathbb Z\delta}
\end{tikzcd}}

(the isomorphism
\diagram{\begin{tikzcd}[column sep=large,row sep=small]
&&\inserted{\underline\Sigma_{\rho,\sigma}^{\mathrm{eff}+0}}\\
\mathbb G_a\arrow[r,phantom,"\simeq"]&\underline Z_{\rho,\sigma}^{+0}\arrow[ur,"\inserted{\simeq}"]&
\deleted{\unsure{\mathbb G_a}}\arrow[l,phantom,"\deleted{\unsure{\xleftarrow{\simeq}}}"']
\end{tikzcd}}
being given by
\diagram{\begin{tikzcd}[column sep=large]
\lambda\arrow[r,mapsto]&(\underline\mu=1,\xi_0\deleted{\unsure{=\eta_0}}=\lambda\delta,\eta_0=\xi_0=\lambda\delta)
\end{tikzcd}}
)

We \deleted{shall have} shall set
\[
(143)\qquad
\unsure{\underline L_0^{\sharp}|L_0^{\sharp}}\overset{\mathrm{dfn}}{=}
\{\lambda\mathbin{.}\lambda_0\mid\lambda\in\Gamma\mathbb G_a\}\subset\underline M
\]
and we define

\((144)\)
\diagram{\begin{tikzcd}[column sep=large]
\underline L_0^{\sharp}\arrow[r]&\underline Z_{\rho,\sigma}^{\unsure{+0}}\arrow[r,"\simeq","\bigl("{description,pos=0,anchor=east}]&
\underline\Sigma_{\rho,\sigma}^{\mathrm{eff}\unsure{+0}}\bigr)\\
\underset{\prime\prime\lambda\mathbin{.}\lambda_0}{\ell}\arrow[r,mapsto]&
(\operatorname{int}(\ell),1)\arrow[r,mapsto]&
\begin{gathered}
(\deleted{\unsure{\mu}}1,(1-\rho_{\underline M})(\ell),(1-\sigma_{\underline M})(\ell))\\
=(1,-\lambda\delta,-\lambda\delta)
\end{gathered}
\end{tikzcd}}

\scanpage{76}

thus we have
\[
(145)\qquad\underline L_0^{\sharp}\simeq\underline Z_{\rho,\sigma}^{+0}
\]
we shall have
\[
(146)\qquad\underline\gamma_{\rho,\sigma}\overset{\mathrm{dfn}}{=}
\underline Z_{\rho,\sigma}/\underline L_0^{\sharp}\simeq\underline{\mathbb O}_{12}^{*}\ .
\]
If we now construct \(G_{\rho,\sigma}\)

\deleted{\((\unsure{147})\)}

as
\[
(147)\qquad\underline G_{\rho,\sigma}
=(\underline Z_{\rho,\sigma}\mathbin{.}\underline{\mathcal G}^{+})/\underline L_0^{\sharp}
\]
we find

\((148)\)
\diagram{\begin{tikzcd}[column sep=large,row sep=large]
1\arrow[r]&\underline L_0^{\sharp}\arrow[r]\arrow[d]&
\underline Z_{\rho,\sigma}\arrow[r]\arrow[d]&
\overbrace{\gamma_{\rho,\sigma}}^{\mathbb O_{12}^{*}}\arrow[r]\arrow[d,phantom,"\rotatebox{90}{$\simeq$}" description]&1\\
1\arrow[r,dash]&\underline{\mathcal G}^{+}\arrow[r]&G_{\rho,\sigma}\arrow[r]&\gamma_{\rho,\sigma}\arrow[r]&1
\end{tikzcd}}
.

Moreover we have a canonical splitting
of \(\underline Z_{\rho,\sigma}\) over \(\gamma_{\rho,\sigma}=\mathbb O_{12}^{*}\)

\((149)\)
\diagram{\begin{tikzcd}[column sep=large]
\mathbb O_{12}^{\unsure{*|\wedge}}\arrow[r]&\underline Z_{\rho,\sigma}\arrow[r,"\simeq"]&\underline\Sigma_{\rho,\sigma}^{\mathrm{eff}}\\
\underline\mu\arrow[r,mapsto]&\unsure{u_{\underline\mu}|\underline\mu}\arrow[r]&
(\underline\mu,\xi_0=\nu''(3\lambda_1),\eta_0=0)
\end{tikzcd}}

\margin{\uline{N.B.} We set
\(G_{\rho,\sigma}^{\sharp}=\underbrace{\mathbb O_{12}^{*}}_{\gamma_{\rho,\sigma}}
\mathbin{.}\underset{\underline M}{\underline{\mathcal G}^{\sharp}}\subset G_{\rho,\sigma}\).}

whence
\[
(150)\qquad
\left\{\begin{aligned}
\underline Z_{\rho,\sigma}&\simeq\mathbb O_{12}^{*}\mathbin{.}\underline L_0^{\sharp}\\
G_{\rho,\sigma}&\simeq\mathbb O_{12}^{*}\mathbin{.}\underline{\mathcal G}^{+}
\end{aligned}\right.
\]
The operation of \(\mathbb O_{12}^{*}\) on \(\underline L_0^{\sharp}\) is
via
\diagram{\begin{tikzcd}[column sep=large]
\mathbb O_{12}^{*}\arrow[r,dash]&\mathbb G_m
\end{tikzcd}}
and homotheties, that
of \(\mathbb O_{12}^{*}\) on \(\underline{\mathcal G}^{+}\) is given by
\[
(151)\quad
\left\{\begin{aligned}
u_{\underline\mu}(\rho)&=\bigl[\deleted{\unsure{3\nu''|\nu_3}}\unsure{6\nu''|6\nu}\lambda_1\bigr]\rho
\quad\text{(}=\lambda_1^{\unsure{6\nu''|6\nu}}\rho\quad\text{in \deleted{more} exponential notation}\text{)}\\
u_{\underline\mu}(\sigma)&=\unsure{\omega^{\nu_2}}\sigma\\
\deleted{\unsure{u_{\underline\mu}}}\qquad u_{\underline\mu}(x)&=\mu x\qquad\text{if }x\in\Gamma\underline M\ .
\end{aligned}\right.
\]